\chapter{Materiais e métodos}\label{chap:materiais}

Este capítulo apresenta os materiais, ferramentas e procedimentos metodológicos que serão utilizados para o desenvolvimento do presente trabalho. São descritas as ferramentas de software empregadas, a metodologia adotada para a implementação e avaliação dos algoritmos de aprendizado de máquina, os procedimentos de validação dos resultados e o cronograma de atividades previsto para a execução do Trabalho de Conclusão de Curso.

\section{Ferramentas}

Para o desenvolvimento deste trabalho, serão utilizadas ferramentas de software amplamente adotadas pela comunidade científica em tarefas de aprendizado de máquina e análise de dados, todas de caráter gratuito e código aberto.

A linguagem de programação utilizada será o \textit{Python}, em sua versão 3, devido à sua ampla adoção em projetos de ciência de dados e à disponibilidade de bibliotecas especializadas. Para o processamento e manipulação dos dados do conjunto IoT-23, será utilizada a biblioteca \textit{Pandas}, que permite a leitura, limpeza e transformação de grandes volumes de dados tabulares de forma eficiente.

A implementação dos algoritmos de aprendizado de máquina supervisionado (Random Forest, Support Vector Machine e Árvore de Decisão) será realizada com auxílio da biblioteca \textit{Scikit-learn}, que oferece implementações consolidadas desses algoritmos, além de funções para divisão dos conjuntos de treino e teste, normalização dos dados e cálculo das métricas de avaliação. Para a implementação da Rede Neural Convolucional (CNN), será utilizada a biblioteca \textit{Keras}, em conjunto com o \textit{TensorFlow} como backend.

A visualização dos resultados, incluindo gráficos comparativos de desempenho entre os algoritmos e matrizes de confusão, será realizada com as bibliotecas \textit{Matplotlib} e \textit{Seaborn}.

Em relação ao ambiente de execução, será utilizado o \textit{Google Colaboratory} (Google Colab), que oferece um ambiente gratuito em nuvem com suporte a notebooks Jupyter e acesso a recursos de processamento gráfico (GPU), o que viabiliza o treinamento do modelo de CNN sem a necessidade de hardware dedicado local.

Por fim, o conjunto de dados utilizado será o IoT-23, descrito na Seção \ref{sec:dataset}, disponibilizado publicamente pelo \textit{Stratosphere Laboratory} da \textit{Czech Technical University}.

\section{Metodologia para o desenvolvimento}

O desenvolvimento deste trabalho será dividido em cinco etapas principais, descritas a seguir.

\textbf{Etapa 1 — Seleção e pré-processamento dos dados.} Serão selecionados os cenários do dataset IoT-23 que contemplem os ataques mais frequentes (PortScan, DDoS e Okiru, entre outros), além das capturas de tráfego benigno, garantindo representatividade das principais categorias presentes no dataset. Em seguida, será realizado o pré-processamento dos dados, incluindo a remoção de valores ausentes ou inconsistentes, a codificação de atributos categóricos e a normalização dos atributos numéricos.

Em relação à seleção de atributos, serão inicialmente utilizadas todas as \textit{features} relevantes disponibilizadas pelo IoT-23, sendo removidos exclusivamente atributos identificadores (como endereços IP e \textit{timestamps} brutos) e quaisquer campos que possam causar vazamento de informação (\textit{data leakage}) em relação ao rótulo da classe. Caso a análise exploratória indique redundância ou baixa relevância entre os atributos remanescentes, técnicas de seleção de atributos (como análise de importância de \textit{features} no Random Forest ou eliminação recursiva) poderão ser empregadas para reduzir a dimensionalidade do conjunto de dados.

Devido à distribuição desigual entre as classes presentes no dataset IoT-23, na qual classes como PortScan e Okiru são amplamente majoritárias em relação a outras categorias de ataque e ao tráfego benigno, será aplicada a técnica de sobreamostragem sintética \textit{SMOTE} (\textit{Synthetic Minority Over-sampling Technique}) sobre o conjunto de treino, gerando exemplos sintéticos das classes minoritárias por meio de interpolação entre amostras vizinhas. Como alternativa ou complemento, poderá ser utilizado o parâmetro \texttt{class\_weight=`balanced'}, disponível nos modelos Random Forest, SVM e Árvore de Decisão do Scikit-learn, que ajusta o peso de cada classe na função de custo de forma inversamente proporcional à sua frequência. A escolha entre as duas abordagens, ou a combinação de ambas, será avaliada empiricamente durante a fase experimental, considerando seu impacto sobre as métricas de recall e F1-Score das classes minoritárias. É importante destacar que o balanceamento será aplicado exclusivamente sobre o conjunto de treino, de modo a não distorcer a distribuição real das classes no conjunto de teste.

\textbf{Etapa 2 — Divisão dos dados e definição da tarefa de classificação.} O conjunto de dados pré-processado será dividido em subconjuntos de treino e teste, na proporção de 70\% e 30\%, respectivamente, utilizando-se a técnica de amostragem estratificada, de modo a preservar a proporção das classes em ambos os subconjuntos. Essa proporção foi escolhida por ser amplamente adotada em trabalhos da área de detecção de intrusão baseada em aprendizado de máquina, conforme observado em \cite{SouzaArima2024} e \cite{Bochie2020}, representando um equilíbrio adequado entre a quantidade de dados disponíveis para o treinamento dos modelos e a representatividade do conjunto reservado para avaliação.

Os experimentos serão conduzidos em duas configurações de classificação, permitindo uma análise mais completa do problema: (i) uma tarefa \textbf{binária}, na qual o objetivo é distinguir entre tráfego benigno e tráfego malicioso, independentemente do tipo específico de ataque; e (ii) uma tarefa \textbf{multiclasse}, na qual cada tipo de ataque presente no dataset (PortScan, DDoS, Okiru, C\&C, entre outros) é tratado como uma classe distinta, além da classe de tráfego benigno. A avaliação na configuração binária permite mensurar a capacidade geral dos modelos em detectar comportamento anômalo, enquanto a configuração multiclasse possibilita avaliar a capacidade de cada modelo em discriminar entre os diferentes tipos de ataque, informação relevante para a tomada de decisão em sistemas de resposta a incidentes.

\textbf{Etapa 3 — Implementação dos modelos.} Serão implementados quatro modelos de aprendizado de máquina: Random Forest, Support Vector Machine (SVM), Árvore de Decisão e uma Rede Neural Convolucional (CNN). A escolha desses algoritmos justifica-se por representarem diferentes famílias de classificadores amplamente empregadas na literatura de detecção de intrusão: métodos baseados em conjuntos de árvores (Random Forest), métodos baseados em uma única estrutura hierárquica de regras (Árvore de Decisão), métodos baseados em margens de separação entre classes (SVM) e métodos de aprendizado profundo (CNN). Essa diversidade permite comparar abordagens com diferentes níveis de complexidade, interpretabilidade e custo computacional, alinhando-se aos algoritmos utilizados no principal trabalho de referência adotado neste estudo \cite{SouzaArima2024}. Os três primeiros modelos serão treinados utilizando diretamente os atributos do tráfego de rede extraídos do dataset, em formato tabular (vetores de características).

Embora as CNNs tenham sido originalmente desenvolvidas para tarefas de visão computacional, sua aplicação à classificação de tráfego de rede vem sendo investigada na literatura de detecção de intrusão, uma vez que a reorganização dos atributos em representações matriciais permite à rede explorar sua capacidade de identificar padrões locais e relações de vizinhança entre características do tráfego, de forma análoga à extração de padrões espaciais em imagens. Para a CNN, os dados serão convertidos em uma representação matricial bidimensional, de forma análoga a uma imagem. Cada amostra de tráfego, originalmente representada por um vetor de $n$ atributos numéricos (por exemplo, duração do fluxo, número de pacotes, número de bytes trocados, portas de origem e destino, entre outros), será reorganizada em uma matriz de dimensões fixas por meio do redimensionamento (\textit{reshape}) do vetor de atributos, com aplicação de \textit{zero-padding} quando necessário, conforme abordagens encontradas na literatura de detecção de intrusão com aprendizado profundo.

\textbf{Etapa 4 — Ajuste de hiperparâmetros.} Para os modelos clássicos (Random Forest, SVM e Árvore de Decisão), será realizada uma busca de hiperparâmetros mais ampla, por meio da técnica de \textit{Grid Search} combinada com validação cruzada estratificada de 5 folds (\textit{5-fold stratified cross-validation}), preservando a proporção das classes em cada fold de forma consistente com a estratégia de amostragem adotada na Etapa 2, explorando, por exemplo, o número de estimadores e a profundidade máxima das árvores no Random Forest, o tipo de kernel e os parâmetros $C$ e $\gamma$ no SVM, e o critério de divisão e a profundidade máxima na Árvore de Decisão.

Para a CNN, dado o elevado custo computacional do treinamento de redes neurais, a busca de hiperparâmetros será mais restrita, limitando-se a um conjunto reduzido de configurações pré-selecionadas (\textit{Random Search} com número limitado de combinações), variando-se principalmente o número e o tamanho dos filtros convolucionais, a taxa de aprendizado e o número de épocas de treinamento, com utilização de \textit{early stopping} para evitar treinamentos excessivamente longos. Nesse caso, a validação cruzada de 5 folds completa não será aplicada à CNN; em seu lugar, será utilizada uma divisão adicional do conjunto de treino em treino e validação (\textit{holdout}), de modo a reduzir o custo computacional do processo de ajuste sem comprometer a confiabilidade da seleção de hiperparâmetros.

\textbf{Etapa 5 — Avaliação e comparação dos modelos.} Por fim, os quatro modelos, já com seus hiperparâmetros otimizados, serão avaliados sobre o conjunto de teste, e seus desempenhos serão comparados com base nas métricas descritas na Seção \ref{sec:validacao}, permitindo identificar qual algoritmo apresenta o melhor equilíbrio entre eficácia na detecção de ataques e custo computacional.

Para verificar se as diferenças de desempenho observadas entre os modelos são estatisticamente significativas, serão utilizadas múltiplas execuções experimentais de cada algoritmo. Para os modelos Random Forest, SVM e Árvore de Decisão, essas execuções corresponderão aos resultados obtidos nos 5 \textit{folds} da validação cruzada empregada na Etapa 4. Para a CNN, uma vez que a etapa de ajuste de hiperparâmetros utiliza \textit{holdout} em vez de validação cruzada completa, serão realizadas 5 execuções independentes do treinamento e avaliação do modelo final, com diferentes sementes aleatórias (\textit{random seeds}), mantendo fixos a arquitetura e os hiperparâmetros selecionados na etapa anterior, de modo a gerar um conjunto de 5 valores de desempenho pareável aos dos demais modelos. Sobre esses conjuntos de resultados será aplicado o teste de \textit{Wilcoxon} (\textit{Wilcoxon signed-rank test}, ou teste dos postos sinalizados de Wilcoxon), teste não paramétrico para amostras pareadas, com nível de significância de 5\% ($\alpha = 0,05$), permitindo comparar par a par o desempenho dos quatro modelos e reforçando o rigor científico da comparação realizada.

\section{Validação}\label{sec:validacao}

A validação dos modelos desenvolvidos será realizada por meio de experimentos computacionais, conduzidos sobre o conjunto de teste obtido na Etapa 2 descrita na seção anterior, o qual não será utilizado em nenhuma etapa de treinamento ou ajuste dos modelos.

Para cada um dos quatro algoritmos implementados, serão calculadas as seguintes métricas de avaliação:

\begin{itemize}
    \item \textbf{Acurácia}: proporção de classificações corretas em relação ao total de amostras avaliadas;
    \item \textbf{Precisão}: proporção de instâncias classificadas como ataque que de fato correspondem a ataques reais;
    \item \textbf{Recall (sensibilidade)}: proporção de ataques reais que foram corretamente identificados pelo modelo;
    \item \textbf{F1-Score}: média harmônica entre precisão e recall, utilizada como métrica de equilíbrio entre essas duas medidas;
    \item \textbf{Tempo de inferência}: tempo médio necessário para que o modelo classifique uma única amostra, calculado pela execução da etapa de predição sobre o conjunto de teste completo, dividindo-se o tempo total de execução pelo número de amostras avaliadas. A medição será realizada em ambiente controlado (Google Colab, sem uso concorrente de GPU/CPU por outros processos), repetindo-se o procedimento por múltiplas execuções e reportando-se a média e o desvio padrão dos tempos obtidos, em milissegundos por amostra. Essa métrica é relevante para avaliar a viabilidade de aplicação dos modelos em dispositivos IoT com recursos computacionais limitados.
\end{itemize}

Adicionalmente, embora não constitua uma métrica de avaliação primária, o tempo de treinamento de cada modelo também será registrado, de modo a complementar a análise de viabilidade computacional das abordagens comparadas.

Para os modelos clássicos (Random Forest, SVM e Árvore de Decisão), será utilizada validação cruzada de 5 folds (\textit{5-fold cross-validation}) durante a etapa de treinamento e ajuste de hiperparâmetros, com o objetivo de reduzir o risco de \textit{overfitting} e obter uma estimativa mais robusta do desempenho de cada modelo. As métricas de acurácia, precisão, recall e F1-Score serão calculadas como a média e o desvio padrão dos valores obtidos nos 5 folds. Para a CNN, conforme descrito na Etapa 4 da metodologia, será utilizada uma divisão em treino e validação (\textit{holdout}) em substituição à validação cruzada completa, devido ao custo computacional do treinamento de redes neurais. Cabe destacar que, para fins de comparação estatística entre os modelos por meio do teste de Wilcoxon (conforme descrito na Etapa 5), a CNN será adicionalmente submetida a 5 execuções independentes com diferentes sementes aleatórias, gerando um conjunto de resultados pareável aos das demais técnicas.

Os resultados obtidos para cada algoritmo serão organizados em tabelas comparativas e representados graficamente por meio de gráficos de barras e matrizes de confusão, facilitando a análise comparativa do desempenho relativo entre os modelos avaliados. A partir dessa comparação, será possível discutir qual ou quais algoritmos apresentam maior adequação para a tarefa de detecção de ataques em redes IoT, considerando tanto a eficácia da classificação quanto o custo computacional associado.

Entre as métricas avaliadas, o F1-Score será adotado como principal critério de comparação entre os modelos, em razão do desbalanceamento das classes presente no dataset IoT-23, que pode distorcer a interpretação isolada da acurácia. As demais métricas — acurácia, precisão, recall e tempo de inferência e de treinamento — serão utilizadas como critérios complementares, permitindo uma análise mais completa do compromisso entre eficácia de classificação e custo computacional de cada algoritmo.

\section{Limitações e ameaças à validade}\label{sec:limitacoes}

Embora a metodologia proposta busque seguir rigor científico e reprodutibilidade, algumas limitações devem ser consideradas na interpretação dos resultados deste trabalho. Primeiramente, os experimentos serão conduzidos sobre um único conjunto de dados (IoT-23), de modo que os resultados obtidos podem não se generalizar diretamente para outros \textit{datasets} ou cenários de rede IoT com características distintas. Em segundo lugar, apesar da aplicação de validação cruzada e de técnicas de balanceamento de classes, não é possível eliminar completamente o risco de \textit{overfitting}, especialmente nos modelos com maior número de hiperparâmetros, como Random Forest e CNN. Em terceiro lugar, a utilização do ambiente gratuito do Google Colaboratory impõe limitações computacionais (cotas de uso de GPU, possível indisponibilidade de recursos), que podem restringir a amplitude da busca de hiperparâmetros, sobretudo para a CNN. Por fim, embora o IoT-23 permaneça um \textit{benchmark} amplamente utilizado, suas capturas foram realizadas entre 2018 e 2019, podendo não refletir integralmente padrões de tráfego e técnicas de ataque mais recentes, o que constitui uma limitação de representatividade temporal a ser considerada na discussão dos resultados.

\section{Cronograma}\label{sec:cronograma}

O cronograma de atividades previsto para a execução deste Trabalho de Conclusão de Curso é apresentado na Tabela \ref{tab:cronograma}. As atividades de revisão bibliográfica não constam no cronograma, uma vez que já foram desenvolvidas durante a etapa de elaboração desta proposta.
No mês de julho, as atividades de seleção e pré-processamento dos dados do conjunto IoT-23 serão iniciadas e concluídas até a primeira semana de agosto, contemplando a limpeza, normalização e balanceamento dos dados. Entre agosto e setembro, será realizada a implementação dos três modelos de aprendizado supervisionado tradicional (Random Forest, SVM e Árvore de Decisão), seguida pela implementação do modelo de Rede Neural Convolucional, que se estenderá até outubro devido à maior complexidade de configuração e treinamento. Em outubro e novembro, será conduzido o ajuste de hiperparâmetros dos quatro modelos, buscando otimizar seu desempenho. Na sequência, entre novembro e dezembro, serão executados os experimentos finais, com a coleta dos resultados e o cálculo das métricas de avaliação descritas na Seção \ref{sec:validacao}, seguidos pela análise e discussão comparativa dos modelos. Por fim, a elaboração do texto final do trabalho ocorrerá em paralelo às últimas etapas experimentais, e a defesa do trabalho está prevista para o mês de dezembro.

\begin{table}[htb!]
\centering
\caption{Cronograma de atividades da Proposta de Trabalho de Conclusão.}
\label{tab:cronograma}
\begin{tabular}{|l|llllll|}
\hline
\rowcolor[HTML]{C0C0C0} 
  & \multicolumn{6}{c|}{Ano}                                        \\ 
\cline{2-7} 
\rowcolor[HTML]{C0C0C0} 
\multirow{-2}{*}{Atividades} & \multicolumn{1}{c|}{Jul} & \multicolumn{1}{c|}{Ago} & \multicolumn{1}{c|}{Set} & \multicolumn{1}{c|}{Out} & \multicolumn{1}{c|}{Nov} & Dez \\ \hline
Seleção e pré-processamento dos dados & \multicolumn{1}{c|}{XXXX} & \multicolumn{1}{c|}{X} & \multicolumn{1}{c|}{} & \multicolumn{1}{c|}{} & \multicolumn{1}{c|}{} & \\
\hline
Implementação dos modelos (RF, SVM, AD) & \multicolumn{1}{c|}{} & \multicolumn{1}{c|}{XXX} & \multicolumn{1}{c|}{X} & \multicolumn{1}{c|}{} & \multicolumn{1}{c|}{} & \\
\hline
Implementação do modelo CNN & \multicolumn{1}{c|}{} & \multicolumn{1}{c|}{} & \multicolumn{1}{c|}{XXX} & \multicolumn{1}{c|}{X} & \multicolumn{1}{c|}{} & \\
\hline
Ajuste de hiperparâmetros & \multicolumn{1}{c|}{} & \multicolumn{1}{c|}{} & \multicolumn{1}{c|}{} & \multicolumn{1}{c|}{XXX} & \multicolumn{1}{c|}{X} & \\
\hline
Execução e coleta dos resultados & \multicolumn{1}{c|}{} & \multicolumn{1}{c|}{} & \multicolumn{1}{c|}{} & \multicolumn{1}{c|}{X} & \multicolumn{1}{c|}{XXX} & \\
\hline
Análise e discussão dos resultados & \multicolumn{1}{c|}{} & \multicolumn{1}{c|}{} & \multicolumn{1}{c|}{} & \multicolumn{1}{c|}{} & \multicolumn{1}{c|}{XX} & XX \\
\hline
Elaboração do texto final & \multicolumn{1}{c|}{} & \multicolumn{1}{c|}{} & \multicolumn{1}{c|}{} & \multicolumn{1}{c|}{} & \multicolumn{1}{c|}{XX} & XX \\
\hline
Banca de defesa & \multicolumn{1}{c|}{} & \multicolumn{1}{c|}{} & \multicolumn{1}{c|}{} & \multicolumn{1}{c|}{} & \multicolumn{1}{c|}{} & X \\
\hline
\end{tabular}
\end{table}